\section{Introduction}

The dominant narrative about redistricting and voter turnout focuses on
the \emph{competitive district hypothesis}: competitive districts produce
higher turnout because voters perceive their vote as more decisive.
This hypothesis is well-supported empirically~\citep{cox1988closed, jacobson2015politics}.

We propose a complementary mechanism: \emph{geographic cohesion}.
Compact districts, because they follow natural geographic features rather
than drawing narrow corridors around specific communities, create constituencies
with stronger shared identity, higher social capital, and more efficient
political mobilization.
These factors should predict higher turnout independent of electoral
competitiveness.

The redistricting connection is direct: algorithmic redistricting produces
more compact districts as a primary objective~\citep{dellaLibera2026b2},
and if compactness predicts turnout, then redistricting method affects
democratic participation.

\paragraph{Mechanism.}
Geographic cohesion operates through three channels:
\begin{enumerate}
  \item \textbf{Shared community identity}: compact districts contain
    communities that are geographically proximate and share local institutions
    (newspapers, civic organizations, schools). Shared identity predicts
    political mobilization~\citep{putnam2000bowling}.
  \item \textbf{Campaign efficiency}: compact districts allow campaigns to
    canvass more efficiently. Door-to-door mobilization produces 5--15 pp
    turnout increases~\citep{gerber2000effects}; compact districts lower
    the cost of achieving this contact.
  \item \textbf{Reduced voter confusion}: irregular district boundaries
    create voter confusion about district membership, which predicts
    reduced turnout at the margin~\citep{mcnulty2009study}.
\end{enumerate}

\paragraph{Identification challenge.}
Compactness and turnout are both endogenous to political conditions.
States with higher political engagement may draw more compact districts
\emph{and} have higher turnout for independent reasons.
We address this with state-cycle fixed effects (absorbing state-specific
engagement and political culture) and use redistricting cycles as
quasi-experimental variation: within a state, compactness changes at
the decade boundary due to redistricting, not due to changes in voter
engagement.
